\paragraph*{} Species interactions structure ecological communities, but predicting when interactions occur and how interactions scale up to community-level outcomes remain central challenges in ecology. Natural systems are characterized by a diverse set of interacting organisms \parencite{Janzen1974Deflowering}; organizing these interaction sets into ecological networks can help us gain insight by linking pairwise interactions to community- and ecosystem-level processes \parencite{lever2014sudden, rohr2014structural}. In the framework of interspecific interaction networks, individual species represent nodes, while the interactions between species represent links. Network approaches have furthered our understanding of a variety of systems reliant on interspecific interactions, including consumer-resource, host-parasite, and mutualistic interactions, and network perspectives have yielded insights on features as specific as the dynamics of singular focal species all the way up to broad-scale patterns of community assembly and stability \parencite{guimaraes2020structure, saravia2022ecological}.

\paragraph*{} One reason network approaches are useful is that they provide a set of descriptive and mathematical tools for characterizing the architecture of species interactions. At the level of individual species, node-level properties can describe the number of interactions a species participates in (degree), the total weight of those interactions (strength), or the position of a species relative to the rest of the network. Centrality, for example, describes how connected a species is to other members of a network and is often used to identify species whose loss may have disproportionate effects on associated communities \parencite{traveset2017plant, maia2019plant, martins2024propagation}. Understanding the ecological and evolutionary processes that give rise to variation in node-level properties is therefore important for understanding how interaction networks are organized and how they may respond to disturbance \parencite{moulatlet2023species, martins2025determinants, zhang2025network}.

\paragraph*{} In addition to node-level characterizations, network approaches can also summarize the overlying architecture of species interactions within a community. Network-level properties include summaries of how links are distributed among species, such as degree distributions and connectance, as well as properties describing the arrangement of links, such as modularity and nestedness \parencite{bascompte2003nested, bascompte2006asymmetric}. These structural properties are of interest because they can influence emergent features of ecological systems, including stability, resilience, robustness to species loss, and the propagation of perturbations through communities \parencite{thebault2010stability, suweis2013emergence, bascompte2022resilience, martins2024propagation}. As a result, understanding the factors that shape network structure can provide insight into both the assembly of ecological communities and the potential responses of those communities to environmental change.

\paragraph*{} The architecture of species interactions emerges from processes operating across multiple scales. Species-level traits \parencite{rafferty2013, ramos2018fruit}, phylogenetic relationships \parencite{eriksson2016evolution}, geographic distributions \parencite{moulatlet2023species}, environmental tolerances \parencite{slatyer2013niche}, local community composition \parencite{saravia2022ecological}, biogeographic context \parencite{sugiura2010, dalsgaard2013, traveset2016, nogales2024review}, and anthropogenic impacts \parencite{teixido2022anthropogenic, cruz2023mutualisms, birkenbach2024land} can all influence which interactions occur and how they are organized. However, the relationship between these factors and interaction structure is not always straightforward. Processes that shape species assemblages may not translate directly into the structure of interactions among those species, and the meaning of network properties may shift between local communities, regional metawebs, and broader macroecological contexts. In this dissertation, I examine how interspecific interactions are organized across these scales in mutualistic, antagonistic, and competitive systems.

\paragraph*{} Like any data collection process, the sampling of ecological networks is imperfect. Even intensive sampling may fail to capture rare interactions or, in some cases, rare interactors \parencite{Young21}. Because biased errors in network construction can alter inferred network structure, incomplete sampling creates challenges for understanding how interactions are organized \parencite{bluthgen2008, Young21}. In chapter two, I address this problem by comparing random-forest regression models informed by combinations of phylogenetic structure, species traits, and latent network structure in their ability to fill in potentially missing interactions in a diverse network of fruiting plants and frugivorous birds in Brazil's Atlantic forest. This chapter highlights the potential importance of latent structural features for predicting mutualistic interactions, and encourages a clear conceptual link between prediction performance metrics and the overall goal of predicting cryptic links.

\paragraph*{} Species’ roles within interaction networks may also depend on processes operating across spatial scales. Highly central species may have disproportionate effects on their associated communities \parencite{traveset2017plant, maia2019plant, martins2024propagation}, but the factors contributing to species' position within interaction networks may differ across biogeographic scales. Identifying the eco-evolutionary processes that give rise to highly central species in local and regional-scale interaction networks is essential for understanding how ecological networks are organized and how they may respond to disturbances \parencite{moulatlet2023species, martins2025determinants, zhang2025network}. In chapter three, I combine a globally distributed dataset of floral visitation networks comprising 30,330 interactions among 6,857 species with large-scale occurrence data to evaluate how species’ geographic and environmental range sizes impact their role in networks across scales. These results suggest that the apparent impacts of geographic range size on species centrality in large-scale metawebs are primarily driven by interaction turnover and the scaling of interaction specificity from local to global networks. When variation in interaction turnover is incorporated, niche breadth exerts the strongest positive effect on plant centrality within both local and global floral visitation networks. 


\paragraph*{} The structure of species interaction networks is shaped by both biotic and abiotic contexts, and the increasing availability of ecological data has permitted examinations into how spatial and environmental gradients influence the structure of these species interaction networks \parencite{tylianakis2017}. In chapter four, I take a macroecological approach to investigate how bipartite interaction networks between free-living hosts and helminth parasites vary across gradients of island area and isolation. I find that relationships between island characteristics and the species richness of both hosts and helminth parasites are broadly consistent with predictions from classical island biogeography theory for free-living species. In contrast, network nestedness showed a more limited response, increasing weakly with island area while failing to respond to either measure of island connectivity. These results suggest that island characteristics shape the richness of host-parasite assemblages more predictably than the structure of interactions within them.

\paragraph*{} Global ecosystems are subject to anthropogenically driven changes at increasing rates, and as a result, so are the mutualistic interaction networks that occur within them. Structural changes in mutualistic networks can be driven by a diverse set of drivers, including shifting climates, urbanization and land use changes, and invasion of non-native species \parencite{traveset2014mutualistic, teixido2022anthropogenic, cruz2023mutualisms}. In chapter five, I characterize the current state of understanding of how these drivers affect the structural properties of two major classes of plant-animal transport mutualisms: seed dispersal and pollination networks \parencite{bronsteinmutualism1}. I also examine how knowledge of these patterns can be used to better understand, predict, and conserve these critical natural systems.

\paragraph*{} Finally, in chapter six, I investigate how species-level performance and interspecific interactions scale up to community assembly outcomes. Single-species vital rates are one possible source of simplified demographic information that may be useful in predicting community outcomes \parencite{ram2019predicting, nestor2023interactions}. Using experimental microbial communities, I test whether demographic parameters derived from monoculture growth experiments can predict multi-strain community outcomes across nutrient environments. Together, these results show that single-species growth parameters can provide useful predictions of community assembly, but that their predictive value depends on environmental contexts.

\paragraph*{} Collectively, these chapters show that the structure and consequences of interspecific interactions are shaped by processes operating across multiple scales. Species traits, phylogenetic relationships, latent network structure, geographic range size, climatic niche breadth, island biogeography, anthropogenic change, and environmental contexts can provide insight into how ecological interactions are organized, but their importance varies across systems and scales. Together, this work contributes to a multiscale view of ecological interactions in which community structure depends not only on whether species interact, but also on how species are embedded within networks, where those networks occur, and how interactions scale up to community-level outcomes.
